% GENERATED FILE. Edit canonical lessons, then run npm run generate:books.
% canonical-source-hash: fnv1a64:75c30a2824f0dd4f
% canonical-lessons: IT-C113-postino, IT-C113-cameriere, IT-C113-cliente, IT-C113-cassa, IT-C113-classe

\chapter{Postman, Waiter, Customer, Till, Class}
\label{ch:it-postman-waiter-customer-till-class}
\hlchaptermodality{113}

\begin{chapteropening}
\textbf{By the end of this chapter:} \emph{I can say a postman, a waiter, a customer, a till and a class.}

Complete the last lesson of chapter 113: I can say a postman, a waiter, a customer, a till and a class.
\end{chapteropening}

\section[il postino]{il postino — a postman}

\label{lesson:IT-C113-postino}



\begin{quote}
Before the new one: say the Italian for an engineer, then the Italian for a firefighter.
\end{quote}

\subsection*{You'll want to know: il postino}
\textbf{il postino} — \textquotedblleft{}a postman\textquotedblright{}.

\begin{cousinweb}
\textbf{postino} — built from \emph{posta}, \textquotedblleft{}the post\textquotedblright{}, with \emph{-ino}: the post's man. EtymDB 2.1 traces \emph{posta} to Latin \emph{postus}, \textquotedblleft{}placed\textquotedblright{}: the post was first a relay station placed along a road.
\end{cousinweb}

\begin{grammarlens}[title={What you've built}]
One more word for this chapter.
\end{grammarlens}

\subsection*{Your turn}
\begin{itemize}
\raggedright
  \item \emph{Say it:} \emph{il postino}
  \item \emph{Say it:} \emph{il postino}, once more
  \item \emph{Say it:} say \emph{il postino}
  \item \emph{Recall:} say the Italian for an engineer, then the Italian for a firefighter, then say \emph{il postino} again
\end{itemize}

\begin{tcolorbox}[breakable,colback=teal!4,colframe=teal!35!black,title={Before you move on}]
What does il postino mean? (\textquotedblleft{}A postman\textquotedblright{}.) Say it once more.
\end{tcolorbox}
\section[il cameriere]{il cameriere — a waiter}

\label{lesson:IT-C113-cameriere}



\begin{quote}
Before the new one: say the Italian for a firefighter, then the Italian for a postman.
\end{quote}

\subsection*{You'll want to know: il cameriere}
\textbf{il cameriere} — \textquotedblleft{}a waiter\textquotedblright{}.

\begin{cousinweb}
\textbf{cameriere} — built from \emph{camera}, \textquotedblleft{}room\textquotedblright{}, with \emph{-iere}, the ending for a man who does a job: first the servant of the rooms. \emph{camera} $\leftarrow$ Latin \emph{camera}, \textquotedblleft{}a vault\textquotedblright{}, from Greek \emph{kamara}. (EtymDB 2.1)
\end{cousinweb}

\begin{grammarlens}[title={What you've built}]
One more word for this chapter.
\end{grammarlens}

\subsection*{Your turn}
\begin{itemize}
\raggedright
  \item \emph{Say it:} \emph{il cameriere}
  \item \emph{Say it:} \emph{il cameriere}, once more
  \item \emph{Say it:} say \emph{il cameriere}
  \item \emph{Recall:} say the Italian for a firefighter, then the Italian for a postman, then say \emph{il cameriere} again
\end{itemize}

\begin{tcolorbox}[breakable,colback=teal!4,colframe=teal!35!black,title={Before you move on}]
What does il cameriere mean? (\textquotedblleft{}A waiter\textquotedblright{}.) Say it once more.
\end{tcolorbox}
\section[il cliente]{il cliente — a customer}

\label{lesson:IT-C113-cliente}



\begin{quote}
Before the new one: say the Italian for a postman, then the Italian for a waiter.
\end{quote}

\subsection*{You'll want to know: il cliente}
\textbf{il cliente} — \textquotedblleft{}a customer\textquotedblright{}.

\begin{cousinweb}
\textbf{cliente} $\leftarrow$ Latin \emph{cliens}, \emph{clientis}, \textquotedblleft{}a client, a follower\textquotedblright{} — a learned borrowing. EtymDB 2.1 has no row for \emph{cliente}; it derives French \emph{client} from \emph{cliens}.
\end{cousinweb}

\begin{grammarlens}[title={What you've built}]
One more word for this chapter.
\end{grammarlens}

\subsection*{Your turn}
\begin{itemize}
\raggedright
  \item \emph{Say it:} \emph{il cliente}
  \item \emph{Say it:} \emph{il cliente}, once more
  \item \emph{Say it:} say \emph{il cliente}
  \item \emph{Recall:} say the Italian for a postman, then the Italian for a waiter, then say \emph{il cliente} again
\end{itemize}

\begin{tcolorbox}[breakable,colback=teal!4,colframe=teal!35!black,title={Before you move on}]
What does il cliente mean? (\textquotedblleft{}A customer\textquotedblright{}.) Say it once more.
\end{tcolorbox}
\section[la cassa]{la cassa — a till, a checkout}

\label{lesson:IT-C113-cassa}



\begin{quote}
Before the new one: say the Italian for a waiter, then the Italian for a customer.
\end{quote}

\subsection*{You'll want to know: la cassa}
\textbf{la cassa} — \textquotedblleft{}a till, a checkout\textquotedblright{}.

\begin{cousinweb}
\textbf{cassa} $\leftarrow$ Latin \emph{capsa}, \textquotedblleft{}a box\textquotedblright{}, from \emph{capere}, \textquotedblleft{}to take, to hold\textquotedblright{} — inherited. (EtymDB 2.1)
\end{cousinweb}

\begin{grammarlens}[title={What you've built}]
One more word for this chapter.
\end{grammarlens}

\subsection*{Your turn}
\begin{itemize}
\raggedright
  \item \emph{Say it:} \emph{la cassa}
  \item \emph{Say it:} \emph{la cassa}, once more
  \item \emph{Say it:} say \emph{la cassa}
  \item \emph{Recall:} say the Italian for a waiter, then the Italian for a customer, then say \emph{la cassa} again
\end{itemize}

\begin{tcolorbox}[breakable,colback=teal!4,colframe=teal!35!black,title={Before you move on}]
What does la cassa mean? (\textquotedblleft{}A till, a checkout\textquotedblright{}.) Say it once more.
\end{tcolorbox}
\section[la classe]{la classe — a class}

\label{lesson:IT-C113-classe}



\begin{quote}
Before the new one: say the Italian for a customer, then the Italian for a till.
\end{quote}

\subsection*{You'll want to know: la classe}
\textbf{la classe} — \textquotedblleft{}a class\textquotedblright{}.

\begin{cousinweb}
\textbf{classe} $\leftarrow$ Latin \emph{classis}, \textquotedblleft{}a division of the people, a class\textquotedblright{} — a learned borrowing: a word passed down by speech would have turned \emph{cl-} into \emph{chi-}. (EtymDB 2.1 records the Latin word.)
\end{cousinweb}

\begin{grammarlens}[title={What you've built}]
One more word for this chapter.
\end{grammarlens}

\subsection*{Your turn}
\begin{itemize}
\raggedright
  \item \emph{Say it:} \emph{la classe}
  \item \emph{Say it:} \emph{la classe}, once more
  \item \emph{Say it:} say \emph{la classe}
  \item \emph{Recall:} say the Italian for a customer, then the Italian for a till, then say \emph{la classe} again
\end{itemize}

\begin{tcolorbox}[breakable,colback=teal!4,colframe=teal!35!black,title={Before you move on}]
What does la classe mean? (\textquotedblleft{}A class\textquotedblright{}.) Say it once more.
\end{tcolorbox}
